\documentclass[11pt,a4paper]{article}
\usepackage[utf8]{inputenc}
\usepackage[T1]{fontenc}
\usepackage[english]{babel}
\usepackage{amsmath,amssymb,amsthm}
\usepackage{geometry}
\geometry{margin=2.5cm}
\usepackage{hyperref}
\usepackage{booktabs}
\usepackage{xcolor}
\usepackage{mdframed}
\usepackage[expansion=false]{microtype}
\usepackage{enumitem}
\usepackage{array}
\usepackage{graphicx}
\usepackage{caption}
\captionsetup{font=small,labelfont=bf,format=plain}
\hypersetup{colorlinks=true, linkcolor=blue!50!black, urlcolor=blue!50!black}

% Epistemic boxes (cohérent avec Papers II/III)
\newmdenv[backgroundcolor=green!10, linecolor=green!50!black, linewidth=1pt,
  innertopmargin=6pt, innerbottommargin=6pt, innerleftmargin=8pt, innerrightmargin=8pt,
  skipabove=8pt, skipbelow=8pt]{rigorousbox}
\newmdenv[backgroundcolor=yellow!15, linecolor=orange!70!black, linewidth=1pt,
  innertopmargin=6pt, innerbottommargin=6pt, innerleftmargin=8pt, innerrightmargin=8pt,
  skipabove=8pt, skipbelow=8pt]{heuristicbox}
\newmdenv[backgroundcolor=blue!8, linecolor=blue!60!black, linewidth=1pt,
  innertopmargin=6pt, innerbottommargin=6pt, innerleftmargin=8pt, innerrightmargin=8pt,
  skipabove=8pt, skipbelow=8pt]{openproblem}
\newmdenv[backgroundcolor=violet!8, linecolor=violet!60!black, linewidth=1pt,
  innertopmargin=6pt, innerbottommargin=6pt, innerleftmargin=8pt, innerrightmargin=8pt,
  skipabove=8pt, skipbelow=8pt]{computational}

\theoremstyle{plain}
\newtheorem{theorem}{Theorem}[section]
\newtheorem{proposition}[theorem]{Proposition}
\newtheorem{corollary}[theorem]{Corollary}
\newtheorem{lemma}[theorem]{Lemma}

\theoremstyle{definition}
\newtheorem{definition}[theorem]{Definition}
\newtheorem{remark}[theorem]{Remark}

\title{\Large Unconditional Primality Certificates\\
for Numbers of the Form $p = 3m(m+1) + 1$:\\[3pt]
\large Arithmetic Filters, Multi-Core Search,\\
and Records to 29\,998 Decimal Digits}
\author{Hassane Bakkaoui\\\small Independent Researcher\\\small bakkahassa@hotmail.com}
\date{2026 --- Paper IV in the parametric prime family series}

\begin{document}
\maketitle

\begin{abstract}
We study the parametric subfamily $p = 3m(m+1) + 1$, $m = 2^a \cdot 3^b - 1$,
$a, b \in \mathbb{N}^*$, from the point of view of unconditional primality
certification via the Pocklington--Lehmer criterion~\cite{pock,brillhart}.

The 3-smooth structure $m+1 = 2^a \cdot 3^b$ provides, for every $(a, b)$, a
fully factored divisor $F = 2^a \cdot 3^{b+1}$ of $p-1$ satisfying
$F > \sqrt{p}$ \emph{unconditionally} (Lemma~\ref{lem:factorisation}), reducing
the certificate to finding witnesses for $q \in \{2, 3\}$ only.

We establish three arithmetic structure theorems:
\begin{itemize}[leftmargin=*]
\item Theorem~\ref{thm:mod6}: $p \equiv 1 \pmod{6}$ for all $(a, b)$.
\item Theorem~\ref{thm:cubic}: every prime $q \equiv 2 \pmod{3}$ satisfies
$q \nmid p$, reducing the sieve by 50\%. Proof: $-3$ is a non-residue mod $q$
when $q \equiv 2 \pmod{3}$.
\item Theorem~\ref{thm:mod7}: $7 \mid p$ if and only if
$(a \bmod 3, b \bmod 6) \in \mathcal{F}_7$, where $\mathcal{F}_7$ is an explicit
set of six pairs, eliminating 33\% of candidates at zero arithmetic cost.
Proof: $\mathrm{ord}_7(2) = 3$, $\mathrm{ord}_7(3) = 6$.
\end{itemize}

Combined, these three results eliminate $\approx 87\%$ of all candidates
\textbf{before any large-number computation}. The PrimeQuest multi-core
algorithm (versions v1--v5) applies this pipeline to produce four unconditional
Pocklington--Lehmer certificates, the largest being
\[
p = 3m(m+1) + 1, \quad m = 2^{19435} \cdot 3^{19173} - 1,
\]
a prime of \textbf{29\,998 decimal digits} found in 3\,h~06\,min on a single
consumer laptop (ARM64, 9 cores).

\smallskip
\textbf{MSC 2020}: 11A41, 11Y11, 11N13.\\
\textbf{Keywords}: primality certificate, Pocklington--Lehmer, centred
hexagonal primes, arithmetic filters, quadratic residues, large prime records.
\end{abstract}

\tableofcontents

\section{Introduction}
\label{sec:intro}

Proving unconditionally that a given integer is prime is a central task of
computational number theory. For generic large primes, the most practical
route is the Pocklington--Lehmer criterion~\cite{pock,brillhart}: exhibit a
fully factored divisor $F > \sqrt{p}$ of $p - 1$ and verify witnesses. The key
difficulty is finding such an $F$.

The subfamily
\[
p = 3m(m+1) + 1, \qquad m = 2^a \cdot 3^b - 1, \quad a, b \in \mathbb{N}^*,
\]
is constructed so that $m + 1 = 2^a \cdot 3^b$ is 3-smooth \emph{by design},
making $F = 2^a \cdot 3^{b+1}$ fully factored and satisfying $F > \sqrt{p}$
for every $(a, b)$ without any condition on $p$.

Beyond the certificate structure, the family carries three independent
arithmetic constraints (Theorems~\ref{thm:mod6}, \ref{thm:cubic},
\ref{thm:mod7}) that collectively remove 87\% of candidates before any
large-number arithmetic. Exploiting all three, the PrimeQuest algorithm has
produced four unconditional certificates up to 29\,998 decimal digits.

\paragraph{Position in the parametric prime family series.}
This paper is the fourth in a series studying the family
$p = km(m+1) + \varepsilon + 2kq$:
\begin{itemize}[leftmargin=*]
\item \textbf{Paper~I}~\cite{paper1}: marginal distribution
$\mathbb{E}[|q_{\min}|] \sim \log m / C_k$, large-scale numerical validation.
\item \textbf{Paper~II}~\cite{paper2}: conditional law
$\mathbb{E}[|q| \mid m] = m/4$ (Conditional Proposition~9.4) and exclusion of
spurious spectral correlations with $\zeta(s)$.
\item \textbf{Paper~III}~\cite{paper3}: $p$-adic distribution of $q_{\min}$
modulo small primes; conditional theorems and a damped-cosine spectral
signature of the lag-$\ell$ autocorrelation.
\item \textbf{Present paper (Paper~IV)}: \emph{unconditional} primality
certification for the 3-smooth subfamily $\varepsilon = +1$, $q = 0$ with
$m + 1 = 2^a 3^b$.
\end{itemize}
The contrast in epistemic status is intentional: Papers~I--III provide
heuristic and conditional analytic results on the full family; Paper~IV
provides \emph{rigorous, verifiable} primality certificates on a structurally
constrained subfamily where the Pocklington method becomes unconditional.

\section{The Pocklington Subfamily}
\label{sec:subfamily}

\subsection{Definition and digit count}

\begin{definition}[Pocklington subfamily]\label{def:subfamily}
For integers $a, b \ge 1$, define
\[
m(a, b) = 2^a \cdot 3^b - 1, \qquad p(a, b) = 3m(m+1) + 1 = 3m^2 + 3m + 1.
\]
The decimal digit count of $p$ is
$D(a, b) = \lfloor 2a \log_{10} 2 + (2b + 1) \log_{10} 3 \rfloor + 1$.
\end{definition}

\subsection{Factorisation of $p-1$ and the certificate structure}

\begin{rigorousbox}
\begin{lemma}[Explicit factorisation]\label{lem:factorisation}
For $p$ as in Definition~\ref{def:subfamily},
\[
p - 1 = 2^a \cdot 3^{b+1} \cdot m, \qquad F := 2^a \cdot 3^{b+1} > \sqrt{p}.
\]
\end{lemma}
\end{rigorousbox}

\begin{proof}
Since $m + 1 = 2^a \cdot 3^b$,
\[
p - 1 = 3m(m+1) = 3m \cdot 2^a \cdot 3^b = 2^a \cdot 3^{b+1} \cdot m.
\]
As $p = 3m^2 + 3m + 1 < 3(m+1)^2$, we get
$\sqrt{p} < \sqrt{3}(m+1) < 3(m+1) = 3 \cdot 2^a \cdot 3^b = 2^a \cdot 3^{b+1} = F$.
\end{proof}

\begin{rigorousbox}
\begin{theorem}[Pocklington--Lehmer certificate~\cite{pock,brillhart}]\label{thm:pock}
Let $p = p(a, b)$ as in Definition~\ref{def:subfamily}. Then $p$ is prime
\emph{if and only if} there exist integers $w_2, w_3$ such that, for $q \in \{2, 3\}$,
\[
w_q^{p-1} \equiv 1 \pmod{p} \qquad \text{and} \qquad
\gcd\!\left(w_q^{(p-1)/q} - 1, \ p\right) = 1.
\]
\end{theorem}
\end{rigorousbox}

\begin{corollary}\label{cor:unconditional}
Theorem~\ref{thm:pock} applies \emph{unconditionally} to every candidate in
the subfamily: by Lemma~\ref{lem:factorisation}, $F > \sqrt{p}$ always holds,
and the certificate requires at most two witnesses ($q \in \{2, 3\}$ only).
\end{corollary}

\begin{remark}\label{rem:witnesses}
In all four certificates produced (Section~\ref{sec:results}), the witnesses
$w_2 = 5$ and $w_3 = 7$ succeeded without exception. See Open
Problem~\ref{op:witnesses}.
\end{remark}

\section{Arithmetic Structure Theorems}
\label{sec:arithmetic}

We establish three independent constraints on the family $p = 3m(m+1) + 1$
with $m = 2^a \cdot 3^b - 1$. Each is proved rigorously and translated into a
concrete computational saving.

\subsection{Congruence modulo 6}

\begin{rigorousbox}
\begin{theorem}[Modulo 6]\label{thm:mod6}
For all $a, b \ge 1$, $p(a, b) \equiv 1 \pmod{6}$.
\end{theorem}
\end{rigorousbox}

\begin{proof}
\textbf{Modulo 2.} Since $m = 2^a \cdot 3^b - 1$ is odd, $m(m+1)$ is even, so
$3m(m+1)$ is even and $p = 3m(m+1) + 1$ is odd.

\textbf{Modulo 3.} We have $3 \mid 3m(m+1)$, hence $p \equiv 1 \pmod{3}$.

Together, $p \equiv 1 \pmod{2}$ and $p \equiv 1 \pmod{3}$ give
$p \equiv 1 \pmod{6}$ by the Chinese Remainder Theorem.
\end{proof}

\begin{corollary}\label{cor:mod6}
No prime $q \in \{2, 3\}$ divides $p(a, b)$, and $p$ lies in the admissible
class $+1$ modulo $6$ for every $(a, b)$.
\end{corollary}

\subsection{Cubic reciprocity sieve}

\begin{rigorousbox}
\begin{theorem}[Cubic reciprocity sieve]\label{thm:cubic}
Let $m \in \mathbb{N}^*$, $p = 3m(m+1) + 1$, and $q \ge 5$ a prime with
$q \equiv 2 \pmod{3}$. Then $q \nmid p$.
\end{theorem}
\end{rigorousbox}

\begin{proof}
The condition $q \mid p$ is equivalent to $3m^2 + 3m + 1 \equiv 0 \pmod{q}$.
Completing the square yields $(6m + 3)^2 \equiv -3 \pmod{q}$, so $-3$ must be
a quadratic residue modulo $q$. By the multiplicativity of the Legendre symbol
and quadratic reciprocity~\cite{ireland-rosen},
\[
\left(\frac{-3}{q}\right)
= \left(\frac{-1}{q}\right)\left(\frac{3}{q}\right)
= (-1)^{\frac{q-1}{2}} \cdot \left(\frac{q}{3}\right) (-1)^{\frac{q-1}{2} \cdot \frac{3-1}{2}}.
\]
A direct evaluation shows $\left(\frac{-3}{q}\right) = +1$ if
$q \equiv 1 \pmod{3}$ and $\left(\frac{-3}{q}\right) = -1$ if
$q \equiv 2 \pmod{3}$. In the latter case, $3m^2 + 3m + 1 \equiv 0 \pmod{q}$
has no solution, so $q \nmid p$.
\end{proof}

\begin{corollary}[Sieve reduction by 50\%]\label{cor:cubic}
In trial division up to a bound $L$, only primes $q \equiv 1 \pmod{3}$ need
be tested. By Dirichlet's theorem, these account for half of all primes in
$[5, L]$. Concretely: of the 78\,498 primes in $[5, 10^6]$, only 39\,231
satisfy $q \equiv 1 \pmod{3}$.
\end{corollary}

\begin{remark}
Theorem~\ref{thm:cubic} holds for the \emph{general} family $p = 3m(m+1) + 1$
with arbitrary $m \in \mathbb{N}^*$, independently of the 3-smooth structure.
\end{remark}

\subsection{Forbidden classes modulo 7}

\begin{rigorousbox}
\begin{theorem}[Forbidden classes modulo 7]\label{thm:mod7}
Let $m = 2^a \cdot 3^b - 1$, $a, b \ge 1$. Define
\[
\mathcal{F}_7 = \bigl\{(0, 2), (0, 3), (1, 0), (1, 1), (2, 4), (2, 5)\bigr\}
\subset \{0, 1, 2\} \times \{0, 1, 2, 3, 4, 5\}.
\]
Then
\[
7 \mid p(a, b) \iff (a \bmod 3, b \bmod 6) \in \mathcal{F}_7.
\]
In particular, exactly $|\mathcal{F}_7|/18 = 1/3$ of all classes are forbidden.
\end{theorem}
\end{rigorousbox}

\begin{proof}
Since $\mathrm{ord}_7(2) = 3$ and $\mathrm{ord}_7(3) = 6$, the value
$m + 1 = 2^a \cdot 3^b \bmod 7$ depends only on $r = a \bmod 3$ and
$s = b \bmod 6$. Setting $t = t(r, s) := 2^r \cdot 3^s \bmod 7$, one has
$m \equiv t - 1 \pmod{7}$ and
\[
p \equiv 3(t-1)t + 1 = 3t^2 - 3t + 1 \pmod{7}.
\]
The condition $7 \mid p$ becomes $3t^2 - 3t + 1 \equiv 0 \pmod{7}$.
Multiplying by $5 \equiv 3^{-1} \pmod{7}$:
\[
t^2 - t - 2 \equiv 0 \pmod{7}, \qquad (t-2)(t+1) \equiv 0 \pmod{7},
\]
so $t \equiv 2$ or $t \equiv 6 \pmod{7}$. Table~\ref{tab:t-rs} lists all
$3 \times 6 = 18$ values of $t(r, s)$; exactly the six pairs in $\mathcal{F}_7$
yield $t \in \{2, 6\}$.
\end{proof}

\begin{table}[ht]
\centering
\caption{Values of $t(r, s) = 2^r \cdot 3^s \bmod 7$. Entries in $\{2, 6\}$
(\textbf{bold}) correspond to $7 \mid p$; these are the six pairs of
$\mathcal{F}_7$.}
\label{tab:t-rs}
\begin{tabular}{c|cccccc}
\toprule
$r \backslash s$ & 0 & 1 & 2 & 3 & 4 & 5 \\
\midrule
0 & 1 & 3 & \textbf{2} & \textbf{6} & 4 & 5 \\
1 & \textbf{2} & \textbf{6} & 4 & 5 & 1 & 3 \\
2 & 4 & 5 & 1 & 3 & \textbf{2} & \textbf{6} \\
\bottomrule
\end{tabular}
\end{table}

\begin{corollary}[Free 33\% elimination]\label{cor:mod7}
Checking $(a \bmod 3, b \bmod 6) \in \mathcal{F}_7$ requires only two modular
reductions and a lookup in a six-element set --- no arithmetic on $D$-digit
numbers. When the search visits pairs $(a, b)$ uniformly over the $18$
residue classes (which holds in the zigzag search of
Section~\ref{sec:algorithm}), exactly $1/3$ of all candidates are rejected
at zero large-number cost.
\end{corollary}

\subsection{Combined elimination rate}

\begin{proposition}[Combined filter efficiency]\label{prop:combined}
Among all candidate pairs $(a, b)$ explored by the zigzag search:
\begin{enumerate}[leftmargin=*,label=(\roman*)]
\item Theorem~\ref{thm:mod7} eliminates 33.3\% (lookup, zero cost).
\item Theorem~\ref{thm:cubic} reduces the sieve to $q \equiv 1 \pmod{3}$;
among the 66.7\% remaining, the sieve eliminates a further $\approx 80\%$.
\item The fraction reaching Miller--Rabin is $\approx 13\%$.
\end{enumerate}
For the 29\,998-digit record (2\,221 pairs tested): 33.3\% by
Theorem~\ref{thm:mod7}, 53.8\% by sieve, 87.1\% total.
\end{proposition}

\section{The PrimeQuest Algorithm}
\label{sec:algorithm}

\subsection{Centre calibration and zigzag}

For target digit count $D$, the search centre is
\[
a_0 = b_0 = \left\lfloor \frac{D}{2 \log_{10} 6} \right\rfloor \approx 0.643 \, D.
\]
The zigzag expands outward: for $\delta = 0, 1, 2, \ldots$ it alternates
$a = a_0 + \delta$ and $a = a_0 - \delta$, enumerating for each $a$ all $b$
with $|D(a, b) - D| \le \tau$ (tolerance $\tau = 10$).

\begin{remark}\label{rem:ratio}
Empirical observation: $b/a = 1.083$ at $D = 19\,999$; $b/a = 0.987$ at
$D = 29\,998$. Since this ratio varies between primes, the symmetric choice
$a_0 = b_0$ minimises the worst-case distance to the prime zone.
\end{remark}

\subsection{Four-stage pipeline}

Each pair $(a, b)$ passes through:
\begin{enumerate}[leftmargin=*]
\item \textbf{Mod-7 filter} (Theorem~\ref{thm:mod7}): reject if
$(a \bmod 3, b \bmod 6) \in \mathcal{F}_7$. Cost: $O(1)$.
\item \textbf{Sieve} (Corollary~\ref{cor:cubic}): for $q \equiv 1 \pmod{3}$,
$q \le L$, compute $p \bmod q$ via small modular exponentiations; reject if
$p \equiv 0$.
\item \textbf{Miller--Rabin}: $k$ rounds on $p(a, b)$ in full precision.
\item \textbf{Pocklington}: if MR passes, search for witnesses $w_2, w_3$ of
Theorem~\ref{thm:pock}.
\end{enumerate}

\subsection{Multi-core execution and checkpoint}

Pairs are distributed across CPU cores via a shared-memory pool. Workers
suppress \texttt{SIGINT}/\texttt{SIGTERM}; checkpoint-resume every 60\,s
enables exact restart after interruption.

\subsection{Algorithm versions}

\begin{table}[ht]
\centering
\caption{Successive versions of PrimeQuest.}
\label{tab:versions}
\begin{tabular}{l l c c}
\toprule
Version & Optimisation added & MR rounds $k$ & Sieve limit $L$ \\
\midrule
v1 & Basic sequential search                       & 25       & $10^6$ \\
v2 & Sieve $\to q \equiv 1 \pmod{3}$ (Thm~\ref{thm:cubic}) & 25       & $10^6$ \\
v3 & Multi-core, checkpoint, centre calibration    & 25       & $10^6$ \\
v4 & Mod-7 filter (Thm~\ref{thm:mod7})             & 25       & $10^6$ \\
v5 & $k: 25 \to 3 + 20^*$; $L: 10^6 \to 10^7$      & $3 + 20$ & $10^7$ \\
\bottomrule
\end{tabular}

\smallskip
\footnotesize
$^*$v5: 3 initial rounds, then 20 confirmation rounds before Pocklington.
\end{table}

\begin{remark}[Correctness of MR round reduction]\label{rem:mr-rounds}
Miller--Rabin is deterministic for primes: every prime passes every round
with certainty. Reducing $k$ from 25 to 3 carries zero risk of missing a true
prime; rare composites passing 3 rounds ($\le 4^{-3} \approx 1.6\%$) are
caught by the 20 confirmation rounds and Pocklington.
\end{remark}

\section{Computational Results}
\label{sec:results}

\subsection{Four certified primes}

\begin{computational}
Table~\ref{tab:certificates} records the four unconditional
Pocklington--Lehmer certificates. In every case:
$w_2 = 5$ for $q = 2$ and $w_3 = 7$ for $q = 3$.
\end{computational}

\begin{table}[ht]
\centering
\caption{Pocklington certificates from the PrimeQuest campaign.}
\label{tab:certificates}
\begin{tabular}{c c c c c c c c}
\toprule
Digits of $p$ & $a$ & $b$ & $b/a$ & MR tests & Elim. & Wall time & Workers \\
\midrule
9\,998   & 6\,212  & 6\,738  & 1.085 & ---    & ---    & ---           & 1 \\
10\,000  & 6\,213  & 6\,740  & 1.085 & ---    & ---    & ---           & 1 \\
19\,999  & 12\,228 & 13\,242 & 1.083 & 1\,050 & 87.9\% & 2\,h~40\,min  & 9 \\
29\,998  & 19\,435 & 19\,173 & 0.987 & 286    & 87.1\% & 3\,h~06\,min  & 9 \\
\bottomrule
\end{tabular}
\end{table}

\subsection{Certificate for the 29\,998-digit prime}

\begin{rigorousbox}
The prime $p = 3m(m+1) + 1$ with $m = 2^{19435} \cdot 3^{19173} - 1$
($14\,999$ digits) has the following \emph{unconditional} certificate:
\begin{align*}
p - 1 &= F \cdot m, \qquad F = 2^{19435} \cdot 3^{19174}\ (15\,000 \text{ digits}),
\quad F > \sqrt{p}, \\
w_2 = 5: \quad & 5^{p-1} \equiv 1 \pmod{p}, \quad \gcd(5^{(p-1)/2} - 1, p) = 1, \\
w_3 = 7: \quad & 7^{p-1} \equiv 1 \pmod{p}, \quad \gcd(7^{(p-1)/3} - 1, p) = 1.
\end{align*}
The leading and trailing digits of $p$ are:
\[
p = 1602087359509060348500 \cdots 16930852697403293697.
\]
\end{rigorousbox}

\subsection{Filter statistics for the 29\,998-digit record}

\noindent
\begin{minipage}{\linewidth}
\centering
\captionof{table}{Filter statistics for the 29\,998-digit record.}
\label{tab:filter-stats}
\begin{tabular}{l c c}
\toprule
Stage                              & Pairs eliminated & Fraction \\
\midrule
Mod-7 filter (Theorem~\ref{thm:mod7})    & 740     & 33.3\% \\
Sieve (Corollary~\ref{cor:cubic})        & 1\,195  & 53.8\% \\
Miller--Rabin (composites)               & 285     & 12.8\% \\
\midrule
Total eliminated                         & 2\,220  & 87.1\% \\
Prime found                              & 1       & 0.045\% \\
\bottomrule
\end{tabular}
\end{minipage}

\medskip

\section{Performance and Scaling}
\label{sec:performance}

\subsection{Dominant cost}

The dominant cost is Miller--Rabin: one round on a $D$-digit modulus costs
$O(D^2 \log D \log \log D)$ bit operations (GMP's FFT
multiplication~\cite{gmp}). Sieve and mod-7 checks are negligible by
comparison.

\subsection{Empirical scaling law}

Wall-clock time per MR event (total time $\div$ MR test count), machine: Dell
Inspiron~14 Plus (Snapdragon X Plus ARM64, 9 workers):

\begin{table}[ht]
\centering
\caption{Empirical wall time per Miller--Rabin event.}
\label{tab:scaling}
\begin{tabular}{c c c c}
\toprule
$D$ & Rounds $k$ & $\tau_{\rm MR}$ (wall) & Version \\
\midrule
19\,999 & 25 & $\approx 9$\,s  & v3 \\
29\,998 & 25 & $\approx 39$\,s & v4 \\
29\,998 & 3  & $\approx 5$\,s  & v5 \\
\bottomrule
\end{tabular}
\end{table}

The ratio $39/9 \approx 4.3$ for $D = 30\,000 / 20\,000 = 1.5\times$ is
consistent with empirical scaling $\tau \propto D^{2.5}$
(Figure~\ref{fig:scaling}).

\begin{figure}[ht]
\centering
\includegraphics[width=0.9\textwidth]{fig_scaling.png}
\caption{Empirical wall time $\tau_{\rm MR}$ per Miller--Rabin event vs.\
decimal digit count $D$ (log-log axes). Three empirical points (orange/red/blue)
are plotted with their respective versions. The dashed lines represent the
theoretical scaling $\tau \propto D^{2.5}$ calibrated separately for v4
($k = 25$, red) and v5 ($k = 3$, blue); their vertical separation reflects
the 8.3$\times$ speedup from round reduction (Remark~\ref{rem:mr-rounds}).
The green star is the projection at $D = 50\,000$ digits with $\pm 20\%$
uncertainty (Section~\ref{sec:projection}).}
\label{fig:scaling}
\end{figure}

\paragraph{Reading the figure.} Three points deserve emphasis:

\begin{itemize}[leftmargin=*]
\item \textbf{The slope $D^{2.5}$ is theoretical, not fitted.} It derives from
the FFT-based GMP multiplication complexity~\cite{gmp}: each MR round performs
$O(\log p) \sim O(D)$ modular multiplications, each at cost $O(D \log D
\log \log D)$, giving $O(D^2 \log D \log \log D) \approx O(D^{2.5})$ over the
range tested. With only two distinct values of $D$ (20\,k and 30\,k), no
empirical fit could discriminate exponents in the range $[2.3, 2.7]$.

\item \textbf{The drop from v4 to v5 at $D = 29\,998$ is a protocol change,
not a scaling effect.} It results from reducing the initial Miller--Rabin
rounds from $k = 25$ to $k = 3$ (Remark~\ref{rem:mr-rounds}), which is
zero-risk for primes (deterministic) and corrected by the 20 confirmation
rounds for composites.

\item \textbf{The projection to 50\,000 digits is an extrapolation, not a
measurement.} The $\pm 20\%$ band reflects (i) cache-effect crossovers in
GMP's FFT thresholds, (ii) thermal throttling on consumer hardware over
multi-hour runs, (iii) non-uniformity of MR-event distribution in the search
zigzag.
\end{itemize}

\subsection{Projection: 50\,000-digit record}
\label{sec:projection}

With v5 ($k = 3$) and scaling $D^{2.5}$:
\[
\tau_{\rm MR}(50\,000) \approx 5 \times \left(\frac{50}{30}\right)^{2.5}
\approx 18\ \text{s/event}.
\]
For 400--500 expected MR events: projected wall time \textbf{6--10 hours}
(9 workers, same machine). A search is currently in progress.

\section{Open Problems}
\label{sec:openpb}

\begin{openproblem}
\textbf{Open Problem~7.1 (Generalised arithmetic filters).}
For primes $\ell \equiv 1 \pmod{3}$, $\ell \in \{13, 19, 31, \ldots\}$, derive
the forbidden set
\[
\mathcal{F}_\ell = \bigl\{(a \bmod \mathrm{ord}_\ell(2),
b \bmod \mathrm{ord}_\ell(3)) : \ell \mid p(a, b)\bigr\}.
\]
For $\ell = 13$: $\mathrm{ord}_{13}(2) = 12$, $\mathrm{ord}_{13}(3) = 3$;
the period is 36 and the filter eliminates $2/13 \approx 15\%$ of additional
candidates. Combining filters for several small $\ell$ could push the pre-MR
elimination rate above 90\%.
\end{openproblem}

\begin{openproblem}
\textbf{Open Problem~7.2 (50\,000-digit certificate).}
Produce an unconditional Pocklington certificate for a prime of 50\,000
decimal digits in the subfamily. By Proposition~\ref{prop:combined},
$\approx 87\%$ of candidates are pre-filtered; projected time is 6--10
hours (v5, ARM64, 9 workers).
\end{openproblem}

\begin{openproblem}\label{op:witnesses}
\textbf{Open Problem~7.3 (Universality of the witnesses $w_2 = 5$, $w_3 = 7$).}
In all four certificates, $w_2 = 5$ and $w_3 = 7$ succeeded. Prove (or
disprove) that 5 and 7 are valid Pocklington witnesses for every prime of
the subfamily.
\end{openproblem}

\section{Conclusion}
\label{sec:conclusion}

We have given a self-contained treatment of unconditional primality
certificates for the 3-smooth subfamily $p = 3m(m+1) + 1$,
$m = 2^a \cdot 3^b - 1$.

Three elementary but effective arithmetic theorems --- the mod-6 identity
(Theorem~\ref{thm:mod6}), the cubic reciprocity sieve
(Theorem~\ref{thm:cubic}), and the mod-7 forbidden-class filter
(Theorem~\ref{thm:mod7}) --- together eliminate 87\% of candidates before
any large-number computation. Their proofs rest on standard tools: the
Chinese Remainder Theorem, the non-residuacity of $-3$ modulo primes
$q \equiv 2 \pmod{3}$, and the multiplicative orders $\mathrm{ord}_7(2) = 3$,
$\mathrm{ord}_7(3) = 6$.

Applied via the PrimeQuest multi-core algorithm on a consumer laptop, this
framework produced four unconditional certificates, culminating in a prime
of 29\,998 decimal digits (286 MR events, 3\,h~06\,min). A fifth search
targeting 50\,000 digits is in progress.

\paragraph{Place in the series.} As Paper~IV of the parametric prime family
series~\cite{paper1,paper2,paper3}, this work complements the analytic and
distributional results of Papers~I--III with a rigorous, computational layer:
where the earlier papers characterise the family heuristically or
conditionally on Bateman--Horn-type assumptions, the Pocklington subfamily
admits unconditional certification. The contrast highlights the value of
identifying structurally distinguished sub-classes within an otherwise
analytically intractable family.

\section*{Acknowledgements}

The author thanks the developers of \texttt{gmpy2} and GMP for making
large-precision arithmetic accessible to independent researchers.
Computations were performed on a Dell Inspiron~14 Plus 7441 (Snapdragon X Plus
ARM64, 10 Oryon cores, LPDDR5X memory).

\begin{thebibliography}{10}

\bibitem{paper1} H.~Bakkaoui, \textit{A Parametric Family of Primes via
Quadratic Forms: Heuristic Distribution of the Minimal Offset, Large-Scale
Numerical Validation, and Statistical Control of Spurious Spectral
Correlations}, Working paper (Paper~I), 2026.

\bibitem{paper2} H.~Bakkaoui, \textit{Parametric Prime Families via Quadratic
Forms: Layer Invariants, Conditional Proofs, and Connections to the Riemann
Zeta Function}, Working paper (Paper~II, revised), 2026.

\bibitem{paper3} H.~Bakkaoui, \textit{$p$-adic Distribution of the Minimal
Offset $q_{\min}$ in Parametric Prime Families: Bateman--Horn Predictions,
Conditional Theorems, and an Anti-Persistence Sign-Reversal Phenomenon},
Working paper (Paper~III), 2026.

\bibitem{aks} M.~Agrawal, N.~Kayal, and N.~Saxena, \textit{PRIMES is in P},
Ann.\ of Math.\ \textbf{160} (2004), 781--793.

\bibitem{gmp} T.~Granlund and the GMP development team, \textit{GNU Multiple
Precision Arithmetic Library}, version~6.3.0,
\url{https://gmplib.org}, 2023.

\bibitem{brillhart} J.~Brillhart, D.~H.~Lehmer, and J.~L.~Selfridge,
\textit{New primality criteria and factorizations of $2^m \pm 1$}, Math.\
Comp.\ \textbf{29} (1975), 620--647.

\bibitem{pock} H.~C.~Pocklington, \textit{The determination of the prime or
composite nature of large numbers by Fermat's theorem}, Proc.\ Cambridge
Philos.\ Soc.\ \textbf{18} (1914--1916), 29--30.

\bibitem{ireland-rosen} K.~Ireland and M.~Rosen, \textit{A Classical
Introduction to Modern Number Theory}, 2nd edn., Springer, 1990.

\end{thebibliography}

\end{document}
